\chapter{理想流体动力学}

\intro{理想流体（ideal fluid）是忽略粘性效应的流体模型。虽然完全无粘的流体在自然界中并不存在，但在许多流动问题中（特别是高雷诺数流动），粘性效应局限于边界层等薄层区域，主流区域的流动可近似为理想流体。本章将从欧拉方程出发，导出伯努利方程，并系统地讨论平面势流理论和升力理论。}

\section{欧拉方程}

\subsection{理想流体的动量方程}

对于无粘流体（理想流体），应力张量简化为各向同性的压强项：
\begin{myformula}
\sigma_{ij} = -p\delta_{ij}
\end{myformula}

代入柯西动量方程 \(\rho D\mathbf{u}/Dt = \Nabla\cdot\boldsymbol{\sigma} + \rho\mathbf{g}\)，注意到 \(\Nabla\cdot(-p\mathbf{I}) = -\Nabla p\)，得到欧拉方程：

\begin{myTheorem}{欧拉方程}
理想流体的运动满足：
\begin{myformula}
\frac{\partial \mathbf{u}}{\partial t} + \mathbf{u}\cdot\Nabla\mathbf{u} = -\frac{1}{\rho}\Nabla p + \mathbf{g}
\end{myformula}
或写为张量形式：
\begin{myformula}
\frac{\partial u_i}{\partial t} + u_j\frac{\partial u_i}{\partial x_j} = -\frac{1}{\rho}\frac{\partial p}{\partial x_i} + g_i
\end{myformula}
\end{myTheorem}

欧拉方程最早由欧拉于1755年导出。其物理含义是：流体微团的加速度等于单位质量流体所受的压力梯度力与体积力之和。与连续性方程组成封闭方程组：

\begin{myformula}
\frac{\partial\rho}{\partial t} + \Nabla\cdot(\rho\mathbf{u}) = 0
\end{myformula}

\begin{myformula}
\rho\left(\frac{\partial \mathbf{u}}{\partial t} + \mathbf{u}\cdot\Nabla\mathbf{u}\right) = -\Nabla p + \rho\mathbf{g}
\end{myformula}

对于不可压缩理想流体（\(\rho = \text{const}\)），简化为：
\begin{myformula}
\Nabla\cdot\mathbf{u} = 0,\quad \frac{\partial \mathbf{u}}{\partial t} + \mathbf{u}\cdot\Nabla\mathbf{u} = -\frac{1}{\rho}\Nabla p + \mathbf{g}
\end{myformula}

共4个方程，4个未知量（\(p\) 和 \(\mathbf{u}\) 的三个分量），封闭可解。

\subsection{欧拉方程与纳维—斯托克斯方程的关系}

欧拉方程是纳维—斯托克斯方程在 \(\mu=0\) 时的特例：

欧拉方程（无粘）：
\[
\frac{D\mathbf{u}}{Dt} = -\frac{1}{\rho}\Nabla p + \mathbf{g}
\]

纳维—斯托克斯方程（有粘，不可压缩）：
\[
\frac{D\mathbf{u}}{Dt} = -\frac{1}{\rho}\Nabla p + \nu\Nabla^2\mathbf{u} + \mathbf{g}
\]

两者的根本区别在于粘性扩散项 \(\nu\Nabla^2\mathbf{u}\)。在理想流体中，由于没有粘性，不能施加无滑移条件（no-slip），只能要求法向速度为零（不可穿透条件）：\(\mathbf{u}\cdot\mathbf{n} = 0\)。这导致理想流体绕流中出现壁面滑移（slip）。

普朗特的边界层理论已证明：在大雷诺数下，无滑移条件的影响限于物面附近的薄边界层内。理想流体理论仍能很好地预测远离物面的主流区流动特征，并为升力计算提供精确结果。

\subsection{欧拉方程的边界条件}

在固体壁面上，理想流体只需满足不可穿透条件：
\begin{myformula}
\mathbf{u}\cdot\mathbf{n} = 0 \quad \text{（在壁面上）}
\end{myformula}

在自由表面上，需满足运动学边界条件（流体质点始终在自由面上）和动力学边界条件（法向应力连续）。

在无穷远处（外流问题），流动趋于均匀来流：
\begin{myformula}
\mathbf{u} \to \mathbf{U}_\infty,\quad p \to p_\infty \quad \text{当 } |\mathbf{x}| \to \infty
\end{myformula}

\section{沿流线的伯努利方程}

\subsection{推导}

伯努利方程是理想流体力学中最重要的积分关系之一。在欧拉方程中，将非线性对流项重写：
\begin{myformula}
\mathbf{u}\cdot\Nabla\mathbf{u} = \Nabla\left(\frac{1}{2}|\mathbf{u}|^2\right) - \mathbf{u}\times(\Nabla\times\mathbf{u})
\end{myformula}

欧拉方程变为兰姆（Lamb）形式：
\begin{myformula}
\frac{\partial\mathbf{u}}{\partial t} + \Nabla\left(\frac{1}{2}|\mathbf{u}|^2\right) - \mathbf{u}\times\boldsymbol{\omega} = -\frac{1}{\rho}\Nabla p + \mathbf{g}
\end{myformula}

考虑以下条件：
\begin{enumerate}
    \item 定常流动：\(\partial\mathbf{u}/\partial t = \mathbf{0}\)
    \item 不可压缩流体：\(\rho = \text{const}\)
    \item 重力为有势力：\(\mathbf{g} = -\Nabla(gz)\)
\end{enumerate}

将兰姆形式的欧拉方程沿流线 \(d\mathbf{r}\) 点乘。沿流线 \(d\mathbf{r}\) 与 \(\mathbf{u}\) 平行，故 \(\mathbf{u}\times\boldsymbol{\omega}\) 垂直于 \(d\mathbf{r}\)，有 \((\mathbf{u}\times\boldsymbol{\omega})\cdot d\mathbf{r} = 0\)。因此：

\[
\Nabla\left(\frac{1}{2}u^2\right)\cdot d\mathbf{r} = -\frac{1}{\rho}\Nabla p \cdot d\mathbf{r} - \Nabla(gz)\cdot d\mathbf{r}
\]

即：
\[
d\left(\frac{1}{2}u^2\right) = -\frac{1}{\rho} dp - g\,dz
\]

积分得：

\begin{myTheorem}{伯努利方程（沿流线）}
对于定常、不可压缩的理想流体，沿同一条流线有：
\begin{myformula}
\frac{p}{\rho} + \frac{u^2}{2} + gz = C(\psi)
\end{myformula}
其中 \(C(\psi)\) 是沿给定流线的常数，不同流线的常数一般不同。
\end{myTheorem}

水头形式为：
\begin{myformula}
\frac{p}{\rho g} + \frac{u^2}{2g} + z = H(\psi)
\end{myformula}

其中 \(p/(\rho g)\) 为压力水头，\(u^2/(2g)\) 为速度水头，\(z\) 为位置水头，三者之和为总水头 \(H\)，沿流线保持不变。

\subsection{无旋流动中的伯努利方程}

如果流动不仅是定常、不可压缩且无旋（\(\boldsymbol{\omega} = \mathbf{0}\)），则 \(\mathbf{u}\times\boldsymbol{\omega} = \mathbf{0}\) 全场成立，伯努利方程在全场均成立：

\begin{myformula}
\frac{p}{\rho} + \frac{u^2}{2} + gz = C \quad \text{（全场常数）}
\end{myformula}

在全场无旋的情况下，常数 \(C\) 在所有流线上相同。这在实际应用中更为方便，不需要追踪流线。

\begin{figure}[H]
\centering
\begin{tikzpicture}[scale=1.3]
    \draw[->] (0,0) -- (6,0) node[right] {$x$};
    \draw[->] (0,0) -- (0,4) node[above] {$z$};
    \draw[blue,thick,domain=0.3:5.5,smooth,samples=40]
        plot(\x, {3-0.3*(\x-2.5)^2});
    \node[blue] at (5,1.5) {流线};
    \fill (1.5,2.7) circle (0.06) node[above left] {$A$};
    \fill (3.5,2.7) circle (0.06) node[above right] {$B$};
    \draw[dashed] (1.5,2.7) -- (1.5,0);
    \draw[dashed] (3.5,2.7) -- (3.5,0);
    \node at (1.5,-0.3) {$z_A$};
    \node at (3.5,-0.3) {$z_B$};
    \node at (1.5,1.0) {$\frac{p_A}{\rho g}$};
    \node at (3.5,1.0) {$\frac{p_B}{\rho g}$};
    \node at (1.5,1.8) {$\frac{u_A^2}{2g}$};
    \node at (3.5,1.8) {$\frac{u_B^2}{2g}$};
    \draw[<->] (1.5,2.7) -- (3.5,2.7);
    \node at (2.5,3.0) {$\frac{p}{\rho}+\frac{u^2}{2}+gz=\text{const}$};
\end{tikzpicture}
\caption{沿流线的伯努利方程示意：总水头守恒}
\label{fig:bernoulli-streamline}
\end{figure}

\section{伯努利方程的应用}

\subsection{皮托管}

皮托管（Pitot tube）由皮托于1732年发明，利用滞止点（stagnation point）原理测量流速。在绕流物体前缘驻点，流速为零，压强为滞止压强 \(p_0\)。沿驻点流线应用伯努利方程（忽略高度差）：

\[
\frac{p_\infty}{\rho} + \frac{U_\infty^2}{2} = \frac{p_0}{\rho} + 0
\]

因此：
\begin{myformula}
U_\infty = \sqrt{\frac{2(p_0 - p_\infty)}{\rho}}
\end{myformula}

实际皮托管（皮托—静压管）的内管测量总压，外管侧壁小孔测量静压。两者之差为动压，经校准后计算流速。广泛应用于航空器（空速管）、风洞实验和工业管道流量测量。

\begin{figure}[H]
\centering
\begin{tikzpicture}[scale=1.3]
    \draw[->] (-0.5,0) -- (5,0) node[right] {$x$};
    \fill[blue!10] (0,-0.5) rectangle (4.5,1.5);
    \draw[thick] (0.8,-0.3) rectangle (1.2,0.8);
    \draw[thick] (1.2,0.25) -- (1.6,0.25);
    \draw[thick,->] (0.2,0.5) -- (0.7,0.5);
    \draw[thick,->] (0.2,0.8) -- (0.7,0.8);
    \draw[thick,->] (0.2,0.2) -- (0.7,0.2);
    \fill[red] (0.8,0.25) circle (0.06);
    \node at (0.8,-0.1) {驻点};
    \node at (2.5,0.9) {静压孔};
    \draw[->] (2.5,0.7) -- (2.0,0.2);
    \node at (4,0.5) {$U_\infty$};
    \node at (3,1.8) {皮托管原理示意};
\end{tikzpicture}
\caption{皮托管原理示意：利用驻点压力与静压之差测速}
\label{fig:pitot}
\end{figure}

\subsection{文丘里管}

文丘里管（Venturi tube）利用伯努利原理测量管道流量，由收缩段、喉部和扩张段组成。

设截面1和截面2（喉部）的面积分别为 \(A_1\) 和 \(A_2\)。由连续性方程和伯努利方程（忽略高度差）：

\[
u_1A_1 = u_2A_2,\quad \frac{p_1}{\rho} + \frac{u_1^2}{2} = \frac{p_2}{\rho} + \frac{u_2^2}{2}
\]

联立解出体积流量：
\begin{myformula}
Q = \frac{A_2}{\sqrt{1 - (A_2/A_1)^2}} \sqrt{\frac{2(p_1-p_2)}{\rho}}
\end{myformula}

引入流量系数 \(C_d\)（通常 \(0.95\sim0.99\)）修正粘性损失。文丘里管因压损小而广泛用于大口径管道流量测量。

\subsection{小孔出流——托里拆利定理}

大容器侧壁有一小孔，液面距小孔高度为 \(h\)。对液面（点1，速度近似为0，压强 \(p_a\)）和小孔出口（点2，压强 \(p_a\)，速度 \(u\)）应用伯努利方程：

\[
\frac{p_a}{\rho} + gh = \frac{p_a}{\rho} + \frac{u^2}{2}
\]

得托里拆利（Torricelli）定理：
\begin{myformula}
u = \sqrt{2gh}
\end{myformula}

即小孔出流速度等于液体从高度 \(h\) 自由下落获得的速度。体积流量为 \(Q = A\sqrt{2gh}\)。考虑射流收缩效应（收缩系数 \(C_c\)）和粘性损失（速度系数 \(C_v\)），实际流量为：

\begin{myformula}
Q = C_d A\sqrt{2gh},\quad C_d = C_c C_v
\end{myformula}

其中 \(C_d\) 为流量系数，锐缘孔口约为 \(0.61\sim0.64\)。

\begin{figure}[H]
\centering
\begin{tikzpicture}[scale=1.2]
    \draw[thick] (0,0) -- (5,0);
    \draw[thick] (0,0) -- (0,4) -- (0.3,4);
    \draw[thick] (4.7,0) -- (5,0) -- (5,4) -- (4.7,4);
    \fill[blue!20] (0,0) rectangle (5,3);
    \node at (2.5,1.5) {液体};
    \draw[->] (5,1.5) -- (6,1.5);
    \node at (5.5,1.8) {$u=\sqrt{2gh}$};
    \draw[<->] (5.5,1.5) -- (5.5,3);
    \node at (5.8,2.25) {$h$};
    \draw[dashed] (0,3) -- (5,3);
    \draw[thick] (0,3.8) -- (0,4.2);
    \draw[thick] (5,3.8) -- (5,4.2);
    \node at (2.5,4.3) {$A_0$};
    \fill[blue] (0,1.2) rectangle (0.3,1.8);
    \node[below] at (0.1,0) {$A$};
\end{tikzpicture}
\caption{容器小孔出流示意图（托里拆利定理）}
\label{fig:torricelli}
\end{figure}

\begin{example}
一大型水槽水面距底部小孔高 \(h=2.5\,\text{m}\)，小孔直径 \(d=3\,\text{cm}\)。设流量系数 \(C_d=0.62\)，求出流速度和流量。

解：
出流理论速度：
\[
u = \sqrt{2gh} = \sqrt{2\times 9.81\times 2.5} = \sqrt{49.05} \approx 7.00\,\text{m/s}
\]

小孔面积：
\[
A = \frac{\pi d^2}{4} = \frac{\pi\times 0.03^2}{4} = 7.07\times 10^{-4}\,\text{m}^2
\]

实际流量：
\[
Q = C_d A\sqrt{2gh} = 0.62\times 7.07\times 10^{-4}\times 7.00 = 3.07\times 10^{-3}\,\text{m}^3/\text{s} \approx 3.1\,\text{L}/\text{s}
\]
\end{example}

\section{无旋流动与速度势理论}

\subsection{速度势与拉普拉斯方程}

如第三章所述，对于无旋流动（\(\Nabla\times\mathbf{u} = \mathbf{0}\)），存在速度势 \(\phi\) 使得 \(\mathbf{u} = \Nabla\phi\)。对于不可压缩流体（\(\Nabla\cdot\mathbf{u} = 0\)），速度势满足拉普拉斯方程：

\begin{myformula}
\Nabla^2\phi = 0
\end{myformula}

拉普拉斯方程是线性方程，基本解可以叠加，这是势流理论应用价值的根本原因。

边界条件：在固体壁面上 \(\partial\phi/\partial n = 0\)（不可穿透条件）；在无穷远处 \(\Nabla\phi \to \mathbf{U}_\infty\)。

\subsection{势流问题的适定性}

拉普拉斯方程是二阶椭圆型偏微分方程，边值问题（Dirichlet或Neumann问题）在适当边界条件下是适定的（well-posed）：解存在、唯一且连续依赖于边界条件。求解方法包括：
\begin{itemize}
    \item 解析方法：分离变量法、保角变换法、基本解叠加法等。
    \item 数值方法：有限差分法、边界元法（BEM）、面元法（panel method）等。
\end{itemize}

面元法在航空航天工业中广泛用于飞行器气动载荷计算。

\section{平面势流与复势}

\subsection{复势的基本理论}

对于二维平面不可压缩无旋流动，速度势 \(\phi(x,y)\) 和流函数 \(\psi(x,y)\) 满足柯西—黎曼条件，可构成解析函数。定义复势（complex potential）：

\begin{mydefinition}{复势}
\begin{myformula}
W(z) = \phi(x,y) + i\psi(x,y)
\end{myformula}
其中 \(z = x + iy\) 为复变量。
\end{mydefinition}

复速度（complex velocity）为：
\begin{myformula}
\frac{dW}{dz} = u - iv
\end{myformula}

其共轭即为复速度矢量：\(\overline{dW/dz} = u + iv\)。速度大小：\(|\mathbf{u}| = |dW/dz|\)。

复势方法将平面势流问题转化为复变函数论问题。任何解析函数 \(W(z)\) 都对应一个可能的平面无旋不可压缩流动——这是势流理论最优雅的特性。

\subsection{基本势流}

\subsubsection{均匀流}

复势：
\begin{myformula}
W(z) = U_\infty z
\end{myformula}

速度分量：\(u = U_\infty\)，\(v = 0\)。若来流与 \(x\) 轴夹角为 \(\alpha\)，则 \(W(z) = U_\infty e^{-i\alpha}z\)。

\subsubsection{源与汇}

位于原点的源（source）的复势：
\begin{myformula}
W(z) = \frac{m}{2\pi}\ln z
\end{myformula}

其中 \(m > 0\) 表示源（流体向外流出），\(m < 0\) 表示汇（流体流向原点）。\(m\) 为体积流量。速度势和流函数：\(\phi = \frac{m}{2\pi}\ln r\)，\(\psi = \frac{m}{2\pi}\theta\)。速度为纯径向：\(u_r = m/(2\pi r)\)，\(u_\theta = 0\)。

\subsubsection{偶极子}

强度为 \(\mu\)、轴线沿 \(x\) 方向的偶极子（doublet）的复势：
\begin{myformula}
W(z) = -\frac{\mu}{2\pi}\frac{1}{z}
\end{myformula}

偶极子为一对等强度但符号相反的源和汇无限接近时的极限，\(\mu = m a\)。速度势和流函数：
\[
\phi = -\frac{\mu}{2\pi}\frac{\cos\theta}{r},\quad
\psi = \frac{\mu}{2\pi}\frac{\sin\theta}{r}
\]

\subsubsection{点涡}

位于原点的强度为 \(\Gamma\)（环量）的点涡（point vortex）的复势：
\begin{myformula}
W(z) = -\frac{i\Gamma}{2\pi}\ln z
\end{myformula}

\(\Gamma > 0\) 对应于逆时针旋转。速度分量：\(u_r = 0\)，\(u_\theta = \Gamma/(2\pi r)\)（纯周向）。速度势和流函数：
\[
\phi = \frac{\Gamma}{2\pi}\theta,\quad
\psi = -\frac{\Gamma}{2\pi}\ln r
\]

点涡在除原点外处处无旋，但绕该奇点的环量为 \(\Gamma\)。

\subsection{叠加原理与圆柱绕流}

将均匀流与偶极子叠加，得到无环量圆柱绕流复势：
\begin{myformula}
W(z) = U_\infty z + \frac{U_\infty a^2}{z}
\end{myformula}

偶极子强度取 \(\mu = 2\pi U_\infty a^2\)。在 \(r = a\) 的圆柱面上径向速度为零，该圆即为绕流物面。

圆柱表面速度分布（极坐标）：
\begin{myformula}
u_\theta(a,\theta) = -2U_\infty\sin\theta
\end{myformula}

圆柱表面压强分布（由全场伯努利方程）：
\begin{myformula}
p(a,\theta) = p_\infty + \frac{1}{2}\rho U_\infty^2\left[1 - 4\sin^2\theta\right]
\end{myformula}

压强系数：
\begin{myformula}
C_p = \frac{p - p_\infty}{\frac{1}{2}\rho U_\infty^2} = 1 - 4\sin^2\theta
\end{myformula}

\begin{figure}[H]
\centering
\begin{tikzpicture}[scale=1.5]
    \draw[thick] (0,0) circle (1.5);
    \node at (0,0) {$a$};
    \draw[->] (-2.5,0) -- (2.5,0) node[right] {$x$};
    \draw[->] (0,-2) -- (0,2) node[above] {$y$};
    \foreach \s in {-2,-1.5,-1,-0.5,0,0.5,1,1.5,2} {
        \draw[blue,->,domain=-2.5:-1.5,smooth,samples=10]
            plot(\x, {\s - 0.3*\s/(\x*\x+\s*\s)});
        \draw[blue,->,domain=1.5:2.5,smooth,samples=10]
            plot(\x, {\s - 0.3*\s/(\x*\x+\s*\s)});
    }
    \foreach \s in {-2,-1.5,-1,-0.5,0,0.5,1,1.5,2} {
        \draw[blue,->,domain=-2:2,smooth,samples=30]
            plot({-1.5 + 0.3*\s*(\x)/(\x*\x+\s*\s)}, \x);
    }
    \fill[red] (0,1.5) circle (0.06) node[above right] {$A$ (驻点)};
    \fill[red] (0,-1.5) circle (0.06) node[below right] {$B$ (驻点)};
    \fill[blue] (1.5,0) circle (0.06) node[right] {$C$};
    \fill[blue] (-1.5,0) circle (0.06) node[left] {$D$};
\end{tikzpicture}
\caption{无环量圆柱绕流的流线图（驻点在前、后缘）}
\label{fig:cylinder-flow}
\end{figure}

\subsection{达朗贝尔佯谬}

从 \(C_p = 1 - 4\sin^2\theta\) 可见，压强分布对 \(x\) 轴和 \(y\) 轴均对称。压强沿圆柱表面积分得到的合力为零——圆柱在理想流体均匀来流中既不受阻力也不受升力。

\begin{myTheorem}{达朗贝尔佯谬}
在不可压缩理想流体的定常均匀绕流中，任意封闭物体所受的合力（阻力和升力）均为零。
\end{myTheorem}

这一结论与日常经验矛盾（物体在流体中运动必受阻力），称为达朗贝尔佯谬（D'Alembert's paradox，1752年）。其根源在于忽略了流体的粘性——实际的边界层和分离流动产生压差阻力。但理想流体理论在升力计算方面是成功的：引入环量即可修正升力为零的结论。

\begin{remark}
达朗贝尔佯谬对流体力学发展影响深远。它促使科学家认识到粘性力的重要性，最终导致纳维—斯托克斯方程（19世纪）和边界层理论（Prandtl，1904年）的建立。认识到理论的失败往往是科学进步的重要驱动力。
\end{remark}

\section{环量与升力}

\subsection{带环量的圆柱绕流}

在无环量圆柱绕流的复势基础上添加点涡（位于圆心），得带环量的圆柱绕流复势：

\begin{myformula}
W(z) = U_\infty\left(z + \frac{a^2}{z}\right) - \frac{i\Gamma}{2\pi}\ln z
\end{myformula}

在圆柱表面 \(r = a\) 上速度分布为：
\begin{myformula}
u_\theta(a,\theta) = -2U_\infty\sin\theta + \frac{\Gamma}{2\pi a}
\end{myformula}

驻点位置由 \(u_\theta = 0\) 确定：\(\sin\theta_s = \Gamma/(4\pi a U_\infty)\)。

当 \(|\Gamma|/(4\pi a U_\infty) < 1\) 时，圆柱面上有两驻点（偏离 \(x\) 轴）；当 \(|\Gamma| = 4\pi a U_\infty\) 时，两驻点合并；当 \(|\Gamma| > 4\pi a U_\infty\) 时，驻点离开圆柱面进入流场。

带环量后，压强分布不对称——上表面压强减小、下表面压强增大（\(\Gamma > 0\)），产生向上的升力。

\subsection{库塔—茹科夫斯基定理}

\begin{myTheorem}{库塔—茹科夫斯基升力定理}
在不可压缩理想流体中，作用在具有环量 \(\Gamma\) 的任意截面柱体上的升力（垂直于来流方向）为：
\begin{myformula}
L = \rho U_\infty \Gamma
\end{myformula}
阻力的理论值仍为零。
\end{myTheorem}

库塔—茹科夫斯基定理（1902--1906年）是空气动力学中最重要的理论结果之一。它揭示了升力的本质机制——环量的存在使绕流不对称，导致压强分布不对称而产生升力。该定理适用于翼型等任意截面二维柱体。

\subsection{库塔条件与环量的确定}

在理想流体框架下，环量 \(\Gamma\) 可以是任意值。但实际翼型绕流有确定环量——库塔（Kutta）条件解决了这一难题。

真实流体中，流体不能绕尖锐后缘流动（否则速度无穷大）。翼型起动过程中产生的起动涡（starting vortex）会调整绕翼型的环量，使得后驻点恰好位于尖锐后缘处，流动在后缘光滑离开。

库塔条件的数学表述：后缘处速度为有限值。对尖后缘翼型，等价于上下翼面后缘流速相等。库塔条件为确定 \(\Gamma\) 提供附加方程，使理想流体理论给出与实验吻合的升力预测。

\begin{example}
某翼型在来流速度 \(U_\infty = 50\,\text{m}/\text{s}\) 的大气中（\(\rho = 1.2\,\text{kg}/\text{m}^3\)），环量 \(\Gamma = 15\,\text{m}^2/\text{s}\)。求单位展长升力。

解：
\[
L = \rho U_\infty \Gamma = 1.2 \times 50 \times 15 = 900\,\text{N}/\text{m}
\]
\end{example}

\subsection{马格努斯效应}

旋转的圆柱体在来流中受到侧向力的现象——马格努斯效应（Magnus effect）。旋转使流体在圆柱一侧加速、另一侧减速，产生环量和相应升力。

在球类运动中非常显著：足球的香蕉球、乒乓球的弧圈球、棒球的曲线球等。马格努斯效应也是弗莱特纳转子船（Flettner rotor ship）的工作原理。

将旋转圆柱简化为均匀流加点涡（在原点）：
\begin{myformula}
W(z) = U_\infty z + \frac{U_\infty a^2}{z} - \frac{i\Gamma}{2\pi}\ln z
\end{myformula}

升力由库塔—茹科夫斯基定理给出。对于旋转圆柱，环量由旋转角速度决定：\(\Gamma \approx 2\pi a^2\Omega\)（粗略估算）。

\begin{figure}[H]
\centering
\begin{tikzpicture}[scale=1.3]
    \draw[thick] (0,0) circle (1.2);
    \draw[->] (0,0) -- (0,1.2) node[midway,left] {$a$};
    \draw[->,thick] (0,1.5) arc[start angle=90, end angle=450, radius=1.5];
    \node at (0.7,2.2) {$\Gamma$};
    \draw[->] (-2.5,0.5) -- (-0.8,0.5);
    \draw[->] (-2.5,0.0) -- (-0.8,0.0);
    \draw[->] (-2.5,-0.5) -- (-0.8,-0.5);
    \node at (-2.2,0.9) {$U_\infty$};
    \draw[->,red,thick] (0,-1.8) -- (0,-2.8) node[below] {$L$ (升力)};
    \foreach \s in {0.5,-0.5} {
        \draw[->,blue,domain=-2.5:-1.5,smooth,samples=10]
            plot(\x, {\s - 0.4*(\x+2)/(\x*\x+\s*\s)});
    }
\end{tikzpicture}
\caption{旋转圆柱绕流——马格努斯效应示意}
\label{fig:magnus}
\end{figure}

\subsection{本章小结}

本章从欧拉方程出发，建立了理想流体动力学的基本框架。导出并详细讨论了伯努利方程及其经典应用（皮托管、文丘里管、托里拆利定理）。研究了无旋流动和速度势理论，利用复势方法系统分析了平面势流的基本解（均匀流、源汇、偶极子、点涡）。通过叠加原理构造了圆柱绕流，讨论了达朗贝尔佯谬及其物理启示。最后，引入环量概念，导出库塔—茹科夫斯基升力定理，讨论了库塔条件和马格努斯效应。理想流体理论虽然在阻力预测方面存在局限性，但在升力计算、外流问题等方面取得了巨大成功，是流体力学理论体系的支柱之一。

\begin{myalgorithm}{圆柱绕流的复势计算与可视化}
\begin{mycode}
import numpy as np
import matplotlib.pyplot as plt

def cylinder_flow_complex(z, U_inf=1.0, a=1.0, Gamma=0.0):
    """
    圆柱绕流复势 W(z) = U*(z + a^2/z) - i*Gamma/(2*pi)*ln(z)
    z: 复平面坐标
    """
    W = U_inf * (z + a**2 / z) - 1j*Gamma/(2*np.pi)*np.log(z)
    return W

def velocity_complex(z, U_inf=1.0, a=1.0, Gamma=0.0):
    """复速度 dW/dz = U*(1 - a^2/z^2) - i*Gamma/(2*pi*z)"""
    dWdz = U_inf * (1 - a**2/z**2) - 1j*Gamma/(2*np.pi*z)
    return dWdz

nx, ny = 200, 150
x = np.linspace(-3, 3, nx)
y = np.linspace(-2, 2, ny)
X, Y = np.meshgrid(x, y)
Z = X + 1j*Y
mask = np.abs(Z) < 1.05
W = cylinder_flow_complex(Z, U_inf=1.0, a=1.0, Gamma=4.0)
psi = W.imag
psi = np.ma.masked_where(mask, psi)
levels = np.linspace(-4, 4, 40)
plt.contour(X, Y, psi, levels=levels, colors='blue')
theta = np.linspace(0, 2*np.pi, 100)
plt.plot(np.cos(theta), np.sin(theta), 'k-', linewidth=2)
plt.axis('equal')
plt.xlabel('$x$')
plt.ylabel('$y$')
plt.title('带环量圆柱绕流流线')
plt.show()
\end{mycode}
\end{myalgorithm}
